\section{Related Work}
\label{sec:related}

\paragraph{Formal theorem proving and language models.}
Neural theorem proving has emerged as a demanding testbed for systematic search
in language models.
LeanDojo \citep{yang2023leandojo} released an open-source toolkit with 98\,K
Lean theorems and the ReProver retrieval-augmented baseline, demonstrating that
tactic prediction can be framed as a retrieval-then-generation problem.
HyperTree Proof Search \citep{lample2022hypertree} applied an AlphaZero-style
MCTS policy to Metamath, achieving 82.6\% success on a standard benchmark.
CoqGym \citep{yang2019coqgym} provided 71\,K Coq proof traces together with the
ASTactic tactic synthesis baseline.
Draft-Sketch-Prove \citep{jiang2023dsp} (ICLR 2023 oral) showed that bridging
informal and formal proofs via sketch generation substantially reduces the
burden on the formal prover.
LEGO-Prover \citep{wang2024legoprover} (ICLR 2024 oral) maintained a growing
skill library that is reused across proofs, improving both success rate and
efficiency.
\gore{} is not a theorem prover, but it exploits the same structural pattern:
all these systems require backtracking search over a tree of tactic applications.
The key distinction is that \gore{} provides a \emph{minimal} substrate
designed specifically to isolate backtracking from the syntactic and
type-theoretic complexity of production proof assistants, while the structural
isomorphism (Section~\ref{sec:isomorphism}) makes existing proof corpora
directly usable as pretraining data.

\paragraph{Reinforcement learning for LLM reasoning.}
GRPO \citep{shao2024deepseekmath} introduced group relative policy optimisation,
which eliminates the critic model by normalising rewards within a group of
sampled outputs, substantially reducing memory overhead.
DAPO \citep{yu2025dapo} extended this with decoupled clip ratios, dynamic
sampling, and token-level policy gradient, improving stability at scale.
Dr.\,GRPO corrects a bias that arises in GRPO when different prompts admit
unequal numbers of valid solutions---a bias that is directly relevant to
\gore{}, where programs with depth $d$ and branching factor $w$ admit $w^d$
valid traces.
A complementary line of work studies \emph{process reward models}
\citep{lightman2023lets}: models trained to score intermediate reasoning steps
rather than final answers.
Training a PRM requires either costly human step-level annotation or a separate
reward model.
\gore{}'s interpreter supplies step-level rewards for free, providing the
cleanest possible setup for studying dense vs.\ sparse reward.

\paragraph{Code-augmented reasoning and task decomposition.}
PAL \citep{gao2023pal} demonstrated that delegating computation to a Python
interpreter---rather than simulating it in natural language---substantially
improves arithmetic and symbolic reasoning.
The same principle underlies Program-of-Thought, Text-to-SQL (Spider
\citep{yu2018spider}, BIRD), and CTE-tree decomposition of complex queries.
In all these settings, an external interpreter verifies correctness, providing
ground-truth reward without human annotation.
\gore{} generalises this idea: where PAL uses Python as the host language and
SQL uses relational algebra, \gore{} uses a minimal logic-programming-style
language whose \callstmt{} nodes correspond exactly to external tool calls in
agent frameworks.
The structural isomorphism between CTE trees and \gore{} programs
(Section~\ref{sec:isomorphism}) means that Spider\,/\,BIRD can serve as a
pretraining source for the \callstmt{} primitive.

\paragraph{Synthetic reasoning environments.}
BabyAI, ARC, and CLUTRR provide controlled tasks for probing reasoning in
language models, but none isolates backtracking search as the target skill.
Tree-of-Thoughts \citep{yao2024tree} and Monte-Carlo Tree Search variants for
LLMs impose a search policy at inference time but do not modify the model's
internal representation of the search.
\gore{} is the first environment where (a) the search tree is a first-class
syntactic object in the model's output, (b) every non-trivial program requires
backtracking, and (c) step-level reward is available without additional
annotation.

\paragraph{Mechanistic interpretability of reasoning.}
The linear representation hypothesis and probing classifiers have revealed that
transformers encode factual attributes in interpretable directions
\citep{lightman2023lets}.
Spectral analyses of residual-stream activations have identified low-rank
structure in the representations of arithmetic and logical operations.
Our Spectral-R1 analysis (Section~\ref{sec:e3}) extends this programme to
backtracking-specific structure: we test whether \textsc{Fork}, \textsc{Cut},
and \textsc{Step} activations occupy distinguishable subspaces and whether
ablating the ``backtrack direction'' causes a measurable performance drop.
Because \gore{} uses synthetic data with no natural-language noise, these
experiments are cleaner than analogous analyses on CoT traces.
